% вариант сделан специально для тренировочного экзамена 2026.04.04, автор: Горелова

\begin{problem}{1}
Найдите значение выражения:
\end{problem}

\samerowans{\begin{problem}{1а (4 балла)}
$\dfrac{4^{6} \cdot 9^{5} + 6^{9} \cdot 120}{8^{4} \cdot 3^{12} - 6^{11}}$
\end{problem}}

\smallskip

\samerowans{\begin{problem}{1б (6 баллов)}
$\dfrac{81^{k+1} - 3^{k+4}}{3^{k+2} \cdot (27^k - 1)}$
\end{problem}}

\smallskip

\samerowans{\begin{problem}{2 (4 балла)}
На координатном луче отмечено несколько точек, координаты которых являются натуральными числами. Известно также, что сумма этих чисел равна~75. Если мы каждую точку переместим на три единичных отрезка, то сумма координат новых точек будет равняться~99. Сколько точек было отмечено на координатном луче?
\end{problem}}

\smallskip

\begin{problem}{3}
Разложите на неразложимые множители (чтобы дальнейшее разложение было невозможно):
\end{problem}

\bigskip

\noindent \textbf{3а (4 балла).} $3a + b(a - 2) - 6 + 4a(2 - a)$

\smallskip

\rowans{1ex}

\bigskip

\noindent \textbf{3б (4 балла).} $-12a^{2}p + 15p^{3} + 8a^{4} - 10a^{2}p^{2}$

\smallskip

\rowans{1ex}

\bigskip

\noindent \textbf{3в* (4 балла).} $3x^{2} + x - 2$

\smallskip

\rowans{1ex}

\smallskip

\begin{problem}{4 (6 баллов)}
Упростите выражение и найдите его значение:
$$(-1{,}5mk)^{3} \cdot \left(-\dfrac{4}{3}m^{2}k\right)^{2} \quad \text{при } m = -1; \quad k = 0{,}5.$$
\end{problem}

\rowans{1ex}

\bigskip

\samerowans{\begin{problem}{5а (4 балла)}
Решите уравнение:
$$8x^{3} - (5 - 3x)(3x + 5) + 11x = (1 + 2x)(4x^{2} + 1 - 2x) - (1 - 9x)(x + 2)$$
\end{problem}}

\bigskip

\begin{problem}{5б (6 баллов)}
Решите уравнение:
$$2|2x - 3| - 5 = 10 - |3 - 2x|$$
\end{problem}


\noindent\grid{34}{12}

\smallskip


\begin{problem}{6 (6 баллов)}
Какое из чисел $A$ и $B$ больше и на сколько:
$$A = \underbrace{666\ldots66}_{25\text{ цифр}} \cdot \underbrace{222\ldots22}_{57\text{ цифр}}
\text{ или } B = \underbrace{333\ldots33}_{25\text{ цифр}} \cdot \underbrace{444\ldots45}_{57\text{ цифр}}$$
\end{problem}

\smallskip

\rowans{1ex}

\smallskip

\begin{problem}{7}
Дана линейная функция $y = 4x + b$. При каком значении~$b$ график этой функции:
\end{problem}

\smallskip

\samerowans{\begin{problem}{7а (4 балла)}
пересекает ось $Ox$ в точке с абсциссой~$-1$?
\end{problem}}

\smallskip

\samerowans{\begin{problem}{7б (4 балла)}
проходит через точку $P(-2;\; 1)$?
\end{problem}}

\smallskip

\samerowans{\begin{problem}{7в (4 балла)}
пересекает ось ординат в точке с ординатой~$5$?
\end{problem}}

\smallskip

\samerowans{\begin{problem}{7г (6 баллов)}
проходит через точку пересечения графиков функций $y = 0{,}5x + 1$ и $y = x - 1$?
\end{problem}}

\smallskip

\begin{problem}{8 (12 баллов)}
Напишите уравнение линейной функции $y = kx + b$, график которой параллелен графику функции $y = x^8 - (x^4 - 1)(x^4 + 1) - 0{,}5(5 - 4x)$ и проходит через начало координат. Постройте этот график.
\end{problem}

\smallskip

\noindent\grid{34}{24}

\smallskip

\begin{problem}{9 (12 баллов)}
Два мотоцикла выехали из города в одном направлении: первый~--- в~8~часов утра, а второй~--- в~10~ч.~50~мин. После того как второй мотоцикл догнал первого, они ещё продолжили путь в течение 2,5~часов. В~момент остановки оказалось, что второй мотоцикл обогнал первый на~30~км. В~котором часу второй мотоцикл догнал первый и какое расстояние проехали мотоциклы до этого, если скорость первого мотоцикла равна~$0{,}6$ скорости второго?
\end{problem}

\smallskip

\noindent\grid{34}{18}

\smallskip

\begin{problem}{10 (10 баллов)}
    На конкурсе выступление каждого участника оценивают семь судей. Каждый судья выставляет оценку~--- целое число баллов от~0 до~10 включительно. Известно, что за выступление участника~<<X>> все члены жюри выставили различные оценки. По старой системе оценивания итоговый балл за выступление определяется как среднее арифметическое всех оценок судей. По новой системе оценивания итоговый балл вычисляется следующим образом: отбрасываются наименьшая и наибольшая оценки и подсчитывается среднее арифметическое пяти оставшихся оценок. В~результате оказалось, что разность итоговых баллов, вычисленных по старой и новой системам оценивания, равна~$\dfrac{1}{35}$. Какие оценки могли быть выставлены судьями участнику~<<X>>?
\end{problem}

\bigskip

\rowans{1ex}

\bigskip

\begin{problem}{11 (12 баллов)}
Решите ОДНО из заданий по выбору.

\smallskip

\noindent а) При каких $a$ уравнение$a(x+0{,}5)-6(x+1)=2x+4$ не имеет корней?

\smallskip

\noindent б) Чему равна сумма $\dfrac{1}{2\cdot5}+\dfrac{1}{5\cdot8}+\dots+\dfrac{1}{95\cdot98}+\dfrac{1}{98\cdot101}$?

\smallskip

\noindent в) Можно ли найти такое число $n$, что $n^3-n$ равно 2026?
\end{problem}

\smallskip

\noindent\grid{34}{24}